%! Author = lili
%! Date = 6/1/26

% \chapter{Grenzwertsätze} #2
%\untit{Vorlesung 20. Juni 2025 (V10 2025)}
\untit{Vorlesung 18. Juni 2026 (V09 2026)}


\section*{Lernziele}
\begin{itemize}
    \item Zentraler Grenzwertsatz:
    Aussage formulieren k¨onnen and anwenden k¨onnen
    \item Poissonscher Grenzwertsatz:
    Aussage formulieren k¨onnen and anwenden k¨onnen
    \item Wann wird welche Approximation verwendet?
\end{itemize}

\begin{er}
    \textbf{Satz von de Moivre-Laplace}\label{er:mlp} \txc{(2026)} \\
    Sei $S_n$ die Anzahl der Erfolge in einem n-fach unabhängig wiederholten Bernoulli-Experiment mit
    Erfolgswahrscheinlichkeit $p \in (0,1)$. \\
    Dann gilt f¨ur jede Wahl von Konstanten a, b mit $- \infty < a < b < \infty$
    \[
        \lim_{n \rightarrow \infty} \pe \big( a < \frac{S_n - np}{\sqrt{npq}} \leq b \big) = \frac{1}{\sqrt{2 \pi}} \intab e^{-x^{2} / 2} dx
    \]
    F¨ur große n folgt insbesondere, dass
    \[
        \pe(np - c < S_n \leq np +c ) \approx
        \txc{\int^{c / \sqrt{npq}}_{-c / \sqrt{npq}} \underbrace{ \frac{1}{\sqrt{2 \pi}} e^{-x^{2} / 2}}_{\vfktname} dx =}
        \frac{1}{\sqrt{2 \pi}} \int^{c / \sqrt{npq}}_{-c / \sqrt{npq}} e^{-x^{2} / 2} dx
        = 2 \vfkt (\frac{c}{\sqrt{npq}}) -1
    \]
    \txc{
        Das folgt direkt aus dem allgemeinen Satz, wenn man die Konstanten a und b passend wählt. \\
        Denn die WS für  $np-c < S_n \leq np+c$ ist äquivalent für die von $-c < S_n - np \leq c$.
        Weil $\sqrt{npq} > 0$ (und das ist gegeben, weil $p \in (0,1)$, also
        $p$ ist echt größer als 0 und echt kleiern als 1;
        $n \leq 1$ ist die Anzahl der Versuche;
        $q = 1-p$ und wir wissen ja bereits, dass
        $p$ echt kleiner 1, also
        $q > 0$), kann man dadurch teilen und das hat dann die passende Form, um Moivre-Laplace anzuwenden:
        \[
            \frac{-c}{\sqrt{npq}} < \frac{S_n - np}{\sqrt{npq}} \leq \frac{c}{\sqrt{npq}}
        \]
    Und die Gleichung ergibt sich aus
    \[
        \vfkt \left(\frac{c}{\sqrt{npq}}\right) - \vfkt \left(-\frac{c}{\sqrt{npq}}\right) = \vfkt \left(\frac{c}{\sqrt{npq}}\right) - \left(1 - \vfkt \left(\frac{c}{\sqrt{npq}}\right)\right) = 2\vfkt \left(\frac{c}{\sqrt{npq}}\right) - 1
    \]
    }
   \begin{itemize}
       \item  Ist $c > 0$ nicht von n abh¨angig, so geht die Ws. mit $n \rightarrow \infty$ gegen Null
       \txc{(c ist eine feste Konstante, hängt nicht von nn
       n ab)}
       \item Ist $ 0 < c << \sqrt{n}$ , so geht die Ws. mit $n \rightarrow \infty$ immer noch gegen Null
       \txc{(c wächst, aber langsamer als $\sqrt{n}$)}
       \item Relevanter Beitrag, d.h., Ws. gr¨oßer als null, nur f¨ur $c \approx \sqrt{n}$ (oder gr¨oßer)
       \txc{(c wächst in derselben Größenordnung wie $\sqrt{n}$))}
   \end{itemize}
    \txc{
        Das erklärt, wie sich die Wahrscheinlichkeit $\mathbb{P}(np-c < S_n \leq np+c) \approx 2\Phi\left(\frac{c}{\sqrt{npq}}\right) - 1$
        verhält, je nachdem wie schnell c mit n wächst (oder eben nicht wächst).
    }
\end{er}

\komm{
Zu \nameref{er:mlp}: Die Liste ist relevant, denn ist npq zu klein, ist die Normalapproximation schlecht, und man sollte lieber exakt mit der Binomialverteilung rechnen (oder mit der Poisson-Approximation, falls
    p sehr klein ist). % TODO notizen.txt
}

\begin{beo}
    \textbf{Vergleich mit dem Gesetz der großen Zahlen} \txc{(2026)}
    \begin{align*}
        \pe(|\frac{1}{n} S_n - \E{X_1} \geq \epsilon|) & = 1 - \pe (- \epsilon < \frac{1}{n} S_n - p < \epsilon) \\
        & \approx 1 - \pe (- \frac{n \epsilon}{\sqrt{npq}} <  \frac{S_n - np}{\sqrt{npq}} \leq \frac{n \epsilon}{
            \sqrt{npq}})
    \end{align*}
    \txc{
    Gesetz der großen Zahlen:
    \[
        \lim_{n \rightarrow \infty} \pe(\mid \frac{1}{n} S_n - \E{X_1} \mid \geq \epsilon) = 0 \quad \forall \epsilon > 0
    \]
    und Satz von de Moivre-Laplace
    \[
        \limnin \pe( a < \frac{S_n - np}{\sqrt{npq}} \leq b) = \frac{1}{\sqrt{2\pi}} \intab e^{x^2 / 2} \dx \txo{ \ =
            \intab \varphi(x) \dx}
    \]
    }
    Verhalten f¨ur beschrieben durch
    \begin{itemize}
        \item $\epsilon$ der Ordnung 1 \txc{(konstant)}: Gesetz der großen Zahlen
        \item $\epsilon$ der Ordnung $\frac{1}{\sqrt{n}}$ \txc{()}
    \end{itemize}
\end{beo}

\komm{
Bneutz man für Stichprobengröße bestimmen, Konfidenzintervalle, Fehlerabschätzung bei Simulationen / Monte-Carlo-Methoden,
    Qualitätskontrolle / Fehlerraten
}

% Merksatz für die Klausur:
% Frage nach "existiert der Grenzwert / konvergiert es" → GGZ reicht
% Frage nach "wie groß muss nn
%n sein" oder "wie genau ist die Schätzung" oder "gib ein Konfidenzintervall an" → ZGWS ist gefragt, mit der Φ\Phi
%Φ-Tabelle

\begin{satz}
    \textbf{Zentraler Grenzwertsatz} \textit{Verallgemeinerung des Satzes von de Moivre-Laplace} \txc{(2026)}  \\
    Gegeben Zufallsvariablen $X_1, X_2, X_3, ...$ auf einem Wahrscheinlichkeitsraum $\wraum$. \\
    Angenommen $X_1, X_2, X_3, ...$ sind unabhängig und identisch verteilt und $\sigma^2 = \Var(X_i) \in (0, \infty)$ für alle i. \\
    (Beachte $ \Var(X_i)  < \infty \Rightarrow \E{X_i}$ ex.) \\
    Dabei ist $S_n = \sumEN X_i$ und $\mathcal{M} = \E{X_i}$. \\
    Dann gilt:
    \[
        \limnin \pe \left(  a < \frac{S_n - n \mathcal{M}}{\sqrt{n \sigma^2}} \leq b \right) = \vfkt(b) - \vfkt(a)
    \]
\end{satz}

\begin{satz}
    \textbf{Zentraler Grenzwertsatz für Bernoulli-Variablen} \txc{(2026)} \\
    Sind die Zufallsvariablen $X_1, X_2, X_3, ...$ unabhängige Bernoulli-Variablen, so
    gelten:
\end{satz}



